% ---------------------------------------------------------------------------
% Related work, organised by what each line of work OWNS, so the contribution is
% stated as a gap rather than as a list.
% ---------------------------------------------------------------------------
\section{Related work}
\label{sec:related}

We organise prior work by the claim it already owns, because our contribution is only
meaningful as the complement of those claims.

\paragraph{Choosing an algorithm per task.} SIA selects among training algorithms at the
granularity of a \emph{task}, and demonstrates that the choice matters. That result is
theirs and we do not re-derive it. What a task-level policy cannot express is that units
\emph{within} one task differ in whether a gradient exists at all: by \eqref{eq:irl} the
silent and informative units of a single task require different estimators, and no
task-level assignment can separate them. Our claim is therefore about granularity, and the
comparison that matters is against SIA at matched compute.

\paragraph{Deciding what to evolve.} Co-Harness co-evolves harness and weights and
establishes the success$\to$model / failure$\to$harness assignment. EvoTrainer co-evolves a
policy with its training harness. Both fix the assignment as a rule over outcomes. We keep
their rule as a special case (\eqref{eq:action} constrained to a partition) and relax it,
because a trajectory read by two different consumers need not be spent on only one of them.

\paragraph{Deciding \emph{when} to evolve.} \emph{Learning, Fast and Slow} separates fast
and slow adaptation and reports the stability benefit of a two-timescale schedule. That
bounds what a cadence claim can be worth, and we treat cadence as a control variable rather
than a contribution.

\paragraph{Clustering trajectories.} MEDS reuses layer-wise logits from the forward pass as
lightweight representations of reasoning behaviour and clusters error patterns with HDBSCAN.
We adopt this directly as the learned instantiation of our cluster key, against a derived
key computed from the silence side of \eqref{eq:irl}; the two are a controlled ablation of
\emph{how} to cluster, holding what is done with the clusters fixed.

\paragraph{Evaluation practice.} Recent evaluation work documents how often agent results
fail to survive controls that were available at the time. We adopt its checks as
requirements rather than as discussion: paired tests on shared items, intervals reported on
the \emph{difference} rather than on each arm, matched-permutation controls that hold the
proportion of units in each mode fixed while shuffling which units get which mode, and a
noise floor established from the same checkpoint scored twice. The matched control is the
one that decides our central claim, because a routing method changes both which units are
routed and how many, and only the permutation control separates those.
